\begin{lemma}
Fix \(n>n_0\), put
\[
M=\noderef{TwoBiteNaturalReducedVertexCount}(n),\qquad
K=\noderef{TwoBiteNaturalIndependenceScale}(1+\varepsilon,n),
\qquad
p=\noderef{TwoBiteEdgeProbability}\!\left(\frac12,n\right),
\]
and assume
\(\noderef{ParameterHierarchyEventualComparisons}
(\varepsilon,\varepsilon_1,\varepsilon_2,n_0)\).
For every fixed \(K\)-set \(I\), the probability that
\(\noderef{FiberAndDegreeEvent}\) holds and the small-class closed-pair bound
for this particular \(I\) fails is at most
\[
\exp\!\left(
-\frac{\binom{K}{2}_{\mathbb R}^2}
 {(\log n)M n^{8\varepsilon}}
\right).
\]
\end{lemma}
\begin{proof}
Fix the \(K\)-set \(I\), and write
\[
L=\log n,\qquad
C=\binom{K}{2}_{\mathbb R}.
\]
Let \(\mathcal D_n\) denote \(\noderef{FiberAndDegreeEvent}\).  For a sample
\(\omega\), put
\[
C_I(\omega)=\noderef{TwoBiteClosedPairsInClass}
  \bigl(\omega,I,\noderef{TwoBiteSmallClass}(\omega,\varepsilon,I)\bigr).
\]
We shall also use that the hierarchy places us in the logarithm-positive
range.  Indeed, the first line of
\(\noderef{ParameterHierarchyEventualComparisons}\) gives \(0<M\) and
\(1\le2pM\).  Since \(M>0\), this implies \(p>0\).  With
\[
p=\noderef{TwoBiteEdgeProbability}\!\left(\frac12,n\right)
 =\frac12\sqrt{L/n},
\]
and since \(n>n_0\ge0\), we have \(n>0\), hence \(L/n>0\) and therefore
\[
L=\log n>0. \tag{0}
\]
Unfolding \(\noderef{ClosedPairsBound}\), failure of the fixed-set small
closed-pair bound is the inequality
\[
|C_I(\omega)|>\varepsilon_1K^2. \tag{1}
\]

We first decompose the closed pairs counted by \(C_I\).  Let
\[
P_R=\pi_R(I),\qquad P_B=\pi_B(I),\qquad P_I=P_R\cup P_B
\]
where \(P_I\) is viewed inside the disjoint red-blue base vertex set.  Define
three finite sets of unordered pairs in \(I\).  First,
\(A_{\mathrm{int}}\) is the cardinality of the union
\[
\bigcup_{x\in S_I\cap P_I}\binom{X_x(I)}2.
\]
Second, \(A_{\mathrm{sf}}\) is the number of unordered pairs
\(\{u,w\}\subseteq I\) such that
\[
\pi_R(u)=\pi_R(w)\quad\hbox{or}\quad \pi_B(u)=\pi_B(w).
\]
Third, \(Z_I\) is the number of unordered pairs \(\{u,w\}\subseteq I\) for
which
\[
\pi_R(u)\ne\pi_R(w),\qquad \pi_B(u)\ne\pi_B(w),
\]
and for which there exists an external small witness
\[
x\in S_I\setminus P_I,\qquad u,w\in X_x(I).
\]
Now let \(\{u,w\}\) be a pair counted by \(C_I\).  By unfolding
\(\noderef{TwoBiteClosedPairsInClass}\), there is a witness
\[
x\in S_I,\qquad \{u,w\}\in\binom{X_x(I)}2.
\]
There are three exhaustive possibilities:
\[
\begin{array}{ll}
\text{internal witness:} & x\in P_I,\\
\text{same-fiber pair:} & x\notin P_I\text{ and }
  \pi_R(u)=\pi_R(w)\text{ or }\pi_B(u)=\pi_B(w),\\
\text{external pair:} & x\notin P_I\text{ and the two red and two blue
  projections are distinct.}
\end{array}
\]
In the first case the pair is counted by the union defining
\(A_{\mathrm{int}}\).  In the second case it is counted by
\(A_{\mathrm{sf}}\).  In the third case it is counted by \(Z_I\).  These three
sets of pairs need not be disjoint, because a pair may have more than one
witness, but the preceding case split gives a covering.  Hence the finite
union bound for cardinalities gives
\[
|C_I|\le A_{\mathrm{int}}+A_{\mathrm{sf}}+Z_I. \tag{2}
\]

On \(\mathcal D_n\), each small witness \(x\in S_I\) satisfies
\[
|X_x(I)|\le\noderef{TwoBiteSmallCutoff}(\varepsilon,n)=n^{2\varepsilon}.
\]
There are at most \(|P_R|+|P_B|\le2K\) internal base vertices.  Therefore,
using \(\binom{a}{2}_{\mathbb R}\le a^2/2\) for \(a\ge0\),
\[
A_{\mathrm{int}}
\le 2K\binom{n^{2\varepsilon}}2_{\mathbb R}
\le K n^{4\varepsilon}. \tag{3}
\]
The first small-cutoff comparison in
\(\noderef{ParameterHierarchyEventualComparisons}\), instantiated at this
\(n>n_0\), is
\[
n^{4\varepsilon}K\le\frac{\varepsilon_1}{2}K^2. \tag{4}
\]
Thus \(A_{\mathrm{int}}\le(\varepsilon_1/2)K^2\) on \(\mathcal D_n\).

Still on \(\mathcal D_n\), the fiber-size clauses say that every red and blue
fiber has size within relative error \(\varepsilon_2\) of \(L^2\).  The
hierarchy gives \(0\le\varepsilon_2\le1\), so every fiber meeting \(I\) has at
most \(2L^2\) vertices of \(I\).  For the red fibers, if the fiber sizes inside
\(I\) are \(q_r=|I\cap F_R(r)|\), then
\[
\sum_r q_r=K,\qquad q_r\le2L^2,
\]
and hence
\[
\sum_r\binom{q_r}{2}_{\mathbb R}
\le \sum_r q_r L^2=K L^2.
\]
The same argument for blue fibers gives another \(K L^2\).  Therefore
\[
A_{\mathrm{sf}}\le2K L^2. \tag{5}
\]

We now estimate \(Z_I\) after conditioning on the injection
\(\pi=\noderef{TwoBiteEmbedding}(\omega)\).  This conditioning fixes all
fibers, all values of \(\pi_R,\pi_B\) on \(I\), and hence the two projection
sets \(P_R,P_B\).  By unfolding \(\noderef{TwoBiteSampleWeight}\), the sample
weight of a sample with red graph \(G_R\), blue graph \(G_B\), and injection
\(\pi\) is
\[
\noderef{GnpGraphWeight}(p,G_R)\,
\noderef{GnpGraphWeight}(p,G_B)\,
\noderef{UniformInjectionWeight}(n,M).
\]
After \(\pi\) is fixed, the uniform-injection factor is a constant.  If there
are no injections, the fiber \(\{\omega:\pi_\omega=\pi\}\) has total weight
zero and there is nothing to prove for it.  Otherwise, conditioning on this
fiber divides by the same positive constant on every sample with injection
\(\pi\).  The remaining conditional law is therefore exactly the product of
the red \(G(M,p)\) edge-coordinate law and the blue \(G(M,p)\)
edge-coordinate law.  In particular, by the product definition
\(\noderef{GnpGraphWeight}\), every finite family of distinct red edge
coordinates is independent with probability \(p\) for each edge required to be
present and \(1-p\) for each edge required to be absent; the same statement
holds in the blue graph, and the red and blue edge-coordinate families are
independent of one another.  Edge coordinates that are not mentioned in a
given event may be summed out, because for each such coordinate the two
weights \(p\) and \(1-p\) add to \(1\).

Fix an unordered pair \(\{u,w\}\subseteq I\) with
\[
\pi_R(u)\ne\pi_R(w),\qquad \pi_B(u)\ne\pi_B(w).
\]
For a fixed red external base vertex \(x\notin P_R\) to witness this pair,
both red
edges
\[
\{x,\pi_R(u)\},\qquad \{x,\pi_R(w)\}
\]
must be present in \(G_R\).  These are distinct base-edge coordinates: the two
endpoints \(\pi_R(u)\) and \(\pi_R(w)\) are distinct, and \(x\notin P_R\), so
neither edge is a loop.  Under the conditional \(\noderef{GnpGraphWeight}\)
law their joint probability is \(p^2\).  Intersecting with the additional
condition \(x\in S_I\) can only decrease the probability.  There are at most
\(M\) choices of such red external \(x\), so the union bound gives a red
contribution at most \(Mp^2\).  The same argument in the blue graph gives at
most \(Mp^2\) for blue external witnesses.  A second union bound over the two
colors gives
\[
\Pr(\{u,w\}\text{ contributes to }Z_I\mid\pi)\le2Mp^2. \tag{6}
\]
Using
\[
p=\frac12\sqrt{L/n}
\]
from \(\noderef{TwoBiteEdgeProbability}\), and the rounding comparison
\[
M\le\frac{n}{L^2}
\]
from \(\noderef{ParameterHierarchyEventualComparisons}\), we get
\[
2Mp^2\le 2\cdot\frac{n}{L^2}\cdot\frac{L}{4n}
=\frac1{2L}. \tag{7}
\]
There are at most \(C=\binom{K}{2}_{\mathbb R}\) unordered pairs in \(I\), so
linearity of expectation gives, uniformly in the fixed injection,
\[
\mu_\pi:=\mathbb E(Z_I\mid\pi)\le\frac{C}{2L}. \tag{8}
\]

Set
\[
t=\frac{C}{\sqrt L}.
\]
The expectation-absorption conjunct of
\(\noderef{ParameterHierarchyEventualComparisons}\), instantiated at this
\(n>n_0\) and with the rounded values \(K,M\), is exactly
\[
\frac{C}{2L}+2K L^2+\frac{C}{\sqrt L}
\le \frac{2C}{\sqrt L}. \tag{9}
\]
The final small-pairs comparison in the same package is
\[
\frac{2C}{\sqrt L}\le\frac{\varepsilon_1}{2}K^2. \tag{10}
\]
Consequently, if \(Z_I<\mu_\pi+t\), then by (5), (8), (9), and (10),
\[
A_{\mathrm{sf}}+Z_I
\le 2K L^2+\frac{C}{2L}+\frac{C}{\sqrt L}
\le \frac{\varepsilon_1}{2}K^2. \tag{11}
\]
Combining (2), (3), (4), and (11), no failure of
\(\noderef{ClosedPairsBound}\) can occur on \(\mathcal D_n\).  Therefore
\[
\mathcal D_n\cap\{\text{the fixed-set small bound fails}\}
\subseteq
\left\{Z_I\ge\mu_\pi+t\right\}. \tag{12}
\]

It remains to bound the conditional probability of the event on the right of
(12).  Let
\[
R^\ast=V_R\setminus P_R,\qquad B^\ast=V_B\setminus P_B.
\]
For \(x\in R^\ast\), define the red external adjacency vector
\[
Y_x^R=N_{G_R}(x)\cap P_R,
\]
and for \(y\in B^\ast\), define
\[
Y_y^B=N_{G_B}(y)\cap P_B.
\]
The edge sets used by the distinct vectors \(Y_x^R\) are disjoint, because
all endpoints in \(P_R\) are excluded from \(R^\ast\); the same holds in blue,
and red and blue graphs are independent.  Thus these vectors are independent
under the conditional product law.  More explicitly, if
\[
q_R(A)=p^{|A|}(1-p)^{|P_R|-|A|}\quad(A\subseteq P_R),
\qquad
q_B(B)=p^{|B|}(1-p)^{|P_B|-|B|}\quad(B\subseteq P_B),
\]
then, after all other red and blue edge coordinates are summed out, the
conditional mass of a prescribed family
\((A_x)_{x\in R^\ast}\) and \((B_y)_{y\in B^\ast}\) is
\[
\prod_{x\in R^\ast}q_R(A_x)\prod_{y\in B^\ast}q_B(B_y). \tag{13}
\]
This is precisely the marginal law of the external adjacency vectors.  The
random variable \(Z_I\) is a deterministic function of these vectors.  Indeed,
with \(\pi\) fixed, for \(x\in R^\ast\) one has
\[
X_x(I)=\{u\in I:\pi_R(u)\in Y_x^R\},
\]
and for \(y\in B^\ast\) one has
\[
X_y(I)=\{u\in I:\pi_B(u)\in Y_y^B\}.
\]
Thus whether such an external vertex is in \(S_I\), and which pairs with
distinct projections from \(I\) it contributes to \(Z_I\), are determined by its
own vector and the fixed injection.  All internal vertices and same-fiber
pairs have already been separated off into \(A_{\mathrm{int}}\) and
\(A_{\mathrm{sf}}\), so no other graph coordinates enter \(Z_I\).

The helper \(\noderef{BoundedDifferencesProduct}\) is stated for iid
coordinates.  To match that shape, take one common coordinate space
\[
\Omega=\mathcal P(P_R)\times\mathcal P(P_B)
\]
with product Bernoulli-\(p\) law \(q\) on the two components.  Thus
\[
q(A_R,A_B)=q_R(A_R)q_B(A_B).
\]
It is nonnegative, and
\[
\sum_{(A_R,A_B)\in\Omega}q(A_R,A_B)
=\left(\sum_{A_R\subseteq P_R}q_R(A_R)\right)
 \left(\sum_{A_B\subseteq P_B}q_B(A_B)\right)=1,
\]
because each factor is a finite Bernoulli product distribution.  Enumerate
the vertices of \(R^\ast\) and \(B^\ast\).  For a red external vertex, use the
first component of its \(\Omega\)-coordinate and ignore the second component;
for a blue external vertex, use the second component and ignore the first.
The real external vertices number
\[
|R^\ast|+|B^\ast|\le M+M=2M.
\]
Add dummy coordinates, ignored by \(Z_I\), so that there are exactly
\[
r=2M
\]
coordinates.  Summing over the ignored second components of red coordinates,
the ignored first components of blue coordinates, and all components of dummy
coordinates gives a factor \(1\) in each case.  Therefore the push-forward of
\(q^{\otimes r}\) to the real external adjacency vectors is exactly the
conditional marginal law (13).

Define \(Z:\Omega^r\to\mathbb R\) to be \(Z_I\) computed from these real
external adjacency vectors and to be independent of the ignored and dummy
components.  Written as finite sums, the equality of laws just proved says
\[
\sum_{\eta\in\Omega^r}q^{\otimes r}(\eta)Z(\eta)
=\sum_{\omega:\pi_\omega=\pi}
   \Pr(\omega\mid\pi)\,Z_I(\omega)
=\mathbb E(Z_I\mid\pi).
\]
Thus the expectation of \(Z\) in the product space is exactly
\[
\mu_\pi=\mathbb E(Z_I\mid\pi). \tag{14}
\]

Changing one red coordinate changes only the contribution of the corresponding
external red vertex \(x\).  For a fixed value of all other coordinates, let
\(U\) be the union of external pairs witnessed by all other external vertices.
For a possible value \(A\subseteq P_R\) of the coordinate, let
\[
X_x^A(I)=\{u\in I:\pi_R(u)\in A\}.
\]
The contribution of \(x\) is empty if \(x\) is not small for this value of the
coordinate; otherwise it is a subset of
\[
\binom{X_x^A(I)}2
\]
consisting only of the pairs with two distinct projections.  In the nonempty
case, the definition of the small class
\(\noderef{TwoBiteSmallClass}(\omega,\varepsilon,I)\) gives
\[
|X_x^A(I)|\le\noderef{TwoBiteSmallCutoff}(\varepsilon,n)=n^{2\varepsilon}.
\]
The new value of the coordinate has the same description, with a possibly
different set \(A'\subseteq P_R\).  Put
\[
c=\frac12 n^{4\varepsilon}. \tag{15}
\]
Because \(n>0\), real powers of \(n\) are nonnegative, and hence
\[
\binom{n^{2\varepsilon}}2_{\mathbb R}
\le \frac{(n^{2\varepsilon})^2}{2}
=c,
\]
both the old and new contributed sets have size at most \(c\).  If those two
contributed sets are \(W\) and \(W'\), then the two possible values of \(Z\)
are \(|U\cup W|\) and \(|U\cup W'|\), where \(|W|,|W'|\le c\).  Since
\[
|(U\cup W')\setminus(U\cup W)|\le |W'|\le c,\qquad
|(U\cup W)\setminus(U\cup W')|\le |W|\le c,
\]
their cardinalities differ by at most \(c\).  The same argument applies to one
blue coordinate, and dummy coordinates change \(Z\) by \(0\).  Thus the
Lipschitz hypothesis of \(\noderef{BoundedDifferencesProduct}\) holds with
the exact values
\[
r=2M,\qquad c=\frac12 n^{4\varepsilon},\qquad t=\frac{C}{\sqrt L}.
\]
Here \(c>0\) because \(n>0\) and \(n^{4\varepsilon}>0\), and
\(t\ge0\) because \(C=\binom{K}{2}_{\mathbb R}\ge0\) and \(L>0\) by (0).

Applying \(\noderef{BoundedDifferencesProduct}\) conditionally on \(\pi\)
gives
\[
\Pr(Z_I\ge\mu_\pi+t\mid\pi)
\le
\exp\!\left(-\frac{2t^2}{r c^2}\right).
\]
With the displayed values of \(r,c,t\),
\[
\frac{2t^2}{r c^2}
\;=\;
\frac{2(C^2/L)}
     {(2M)\left(\frac12 n^{4\varepsilon}\right)^2}
\;=\;
\frac{4C^2}{L\,M\,n^{8\varepsilon}}
\;\ge\;
\frac{C^2}{L\,M\,n^{8\varepsilon}}.
\]
Since the exponential function is increasing, replacing the stronger negative
exponent \(-4C^2/(L M n^{8\varepsilon})\) by the weaker
\(-C^2/(L M n^{8\varepsilon})\) gives
\[
\Pr(Z_I\ge\mu_\pi+t\mid\pi)
\le
\exp\!\left(
-\frac{C^2}{L\,M\,n^{8\varepsilon}}
\right). \tag{16}
\]
The right hand side is independent of \(\pi\).  Average (16) over the injection
coordinate.  More explicitly, injections with zero marginal mass contribute
nothing.  For injections with positive marginal mass, the conditional
red-blue product probability was bounded in (16).  Unfolding
\(\noderef{TwoBiteEventProbability}\), the probability of
\(\{Z_I\ge\mu_{\pi_\omega}+t\}\) is the finite sum over injections \(\pi\) of
the injection marginal mass times this conditional probability.  The injection
marginal masses are nonnegative and sum to \(1\), because
\(\noderef{UniformInjectionWeight}(n,M)\) is the uniform injection law.
Therefore the unconditional probability is bounded by the same right hand
side.  Finally, intersecting with \(\mathcal D_n\) can only decrease
probability.  Combining this averaged bound with (12) and substituting
\(C=\binom{K}{2}_{\mathbb R}\) proves
\[
\Pr\!\left(
\mathcal D_n\text{ holds and the fixed-set small closed-pair bound fails}
\right)
\le
\exp\!\left(
-\frac{\binom{K}{2}_{\mathbb R}^2}
 {(\log n)M n^{8\varepsilon}}
\right),
\]
as required.
\end{proof}
